\begin{theorem}\label{thm:kernel-artin-dominant-channel-pole-order}
In the notation of Assumption \ref{prop:kernel-artin-factorization}, let \(G\) be a finite group and let \(u\in(0,1]\). Assume that \(\eta(u)>0\) and that there exist a nontrivial irreducible representation \(\rho_\star\neq \ind\) and an eigenvalue \(\mu_\star\in\mathrm{spec}(M_{K,\rho_\star}(u))\) such that
\begin{enumerate}
  \item \(\abs{\mu_\star}=\eta(u)\);
  \item \(\mu_\star\) occurs in the nontrivial spectrum only in the channel \(\rho_\star\), and there with algebraic multiplicity \(1\); and
  \item every other nontrivial spectral value \(\nu\) satisfies \(\abs{\nu}<\eta(u)\).
\end{enumerate}
Then
\[
\frac{\zeta_{K,G}(z,u)}{L_{K,\ind}(z,u)}
\text{ has a pole of exact order }\dim\rho_\star\text{ at }z=\mu_\star^{-1}.
\]
\end{theorem}

\begin{proof}
By Assumption \ref{prop:kernel-artin-factorization},
\[
\frac{\zeta_{K,G}(z,u)}{L_{K,\ind}(z,u)}
=
\prod_{\rho\in\widehat G\setminus\{\ind\}}L_{K,\rho}(z,u)^{\dim\rho}.
\]
Because \(\mu_\star\) is an eigenvalue of \(M_{K,\rho_\star}(u)\) of algebraic multiplicity \(1\), the polynomial \(\det(I-zM_{K,\rho_\star}(u))\) factors as
\[
\det(I-zM_{K,\rho_\star}(u))
=(1-z\mu_\star)\,g_\star(z),
\]
where \(g_\star\) is a polynomial satisfying \(g_\star(\mu_\star^{-1})\neq 0\). Therefore \(L_{K,\rho_\star}(z,u)=\det(I-zM_{K,\rho_\star}(u))^{-1}\) has a simple pole at \(z=\mu_\star^{-1}\), and
\[
L_{K,\rho_\star}(z,u)^{\dim\rho_\star}
\]
has a pole of order \(\dim\rho_\star\) there.

For every other nontrivial channel \(\rho\neq\rho_\star\), assumptions (2) and (3) imply that \(\mu_\star\notin\mathrm{spec}(M_{K,\rho}(u))\). Hence \(\det(I-zM_{K,\rho}(u))\) is nonzero at \(z=\mu_\star^{-1}\), so \(L_{K,\rho}(z,u)\) is analytic and nonvanishing at that point. Multiplying the factors therefore leaves the pole order unchanged, and the pole order of \(\zeta_{K,G}(z,u)/L_{K,\ind}(z,u)\) at \(z=\mu_\star^{-1}\) is exactly \(\dim\rho_\star\).
\end{proof}

\begin{theorem}\label{thm:kernel-chebotarev-statistical-convergence}
Let \(G\) be a finite abelian group, fix \(u\in(0,1]\), and for all sufficiently large \(n\) define the primitive Chebotarev label distribution
\[
\mu_{n,u}(c):=\frac{B_{n,c}(u)}{B_n(\ind;u)}
\qquad(c\in G).
\]
For any subset \(A\subseteq G\), write
\[
\mu_{n,u}(A):=\sum_{c\in A}\mu_{n,u}(c),
\qquad
\mathrm{unif}_G(A):=\frac{\abs{A}}{\abs{G}},
\]
and let \(\mathrm{unif}_G\) denote the uniform measure on \(G\). Define
\[
\tau(u):=\frac{\Lambda(u)}{\lambda_1(u)}<1,
\qquad
\Lambda(u):=\max\{\abs{\lambda_2(u)},\eta(u)\}.
\]
Then
\[
\|\mu_{n,u}-\mathrm{unif}_G\|_{\mathrm{TV}}
=
O\!\bigl(\tau(u)^n\bigr),
\]
\[
\chi^2\!\bigl(\mu_{n,u}\,\|\,\mathrm{unif}_G\bigr)
=
O\!\bigl(\tau(u)^{2n}\bigr),
\]
\[
D\!\bigl(\mu_{n,u}\,\|\,\mathrm{unif}_G\bigr)
=
O\!\bigl(\tau(u)^{2n}\bigr).
\]
If, in addition, \(\lambda_1(u)>1\) and the square-root spectral bound of Definition \ref{def:kernel-rh} holds, then
\[
\|\mu_{n,u}-\mathrm{unif}_G\|_{\mathrm{TV}}
=
O\!\bigl(\lambda_1(u)^{-n/2}\bigr),
\]
\[
\chi^2\!\bigl(\mu_{n,u}\,\|\,\mathrm{unif}_G\bigr)
=
O\!\bigl(\lambda_1(u)^{-n}\bigr),
\qquad
D\!\bigl(\mu_{n,u}\,\|\,\mathrm{unif}_G\bigr)
=
O\!\bigl(\lambda_1(u)^{-n}\bigr).
\]
Conversely, suppose that there exists a subset \(A\subseteq G\) for which the hypotheses of Theorem \ref{thm:kernel-chebotarev-second-main-term-witness} hold with \(c_A\neq 0\). If \(\lambda_1(u)>1\) and \(\eta(u)>\lambda_1(u)^{1/2}\), then there exists a constant \(c(A)>0\) such that, for infinitely many \(n\),
\[
\bigl|\mu_{n,u}(A)-\mathrm{unif}_G(A)\bigr|
\ge
c(A)\,\frac{\eta(u)^n}{\lambda_1(u)^n}.
\]
In particular, one cannot have
\[
\|\mu_{n,u}-\mathrm{unif}_G\|_{\mathrm{TV}}
=
O\!\bigl(\lambda_1(u)^{-n/2}\bigr)
\]
in that regime.
\end{theorem}

\begin{proof}
Assumption \ref{thm:kernel-chebotarev-exp} gives, uniformly in \(c\in G\),
\[
B_{n,c}(u)
=
\frac{1}{|G|}\frac{\lambda_1(u)^n}{n}
+O\!\left(\frac{\Lambda(u)^n}{n}\right),
\qquad
B_n(\ind;u)
=
\frac{\lambda_1(u)^n}{n}
+O\!\left(\frac{\Lambda(u)^n}{n}\right).
\]
Equivalently,
\[
B_{n,c}(u)
=
\frac{\lambda_1(u)^n}{n}\left(\frac{1}{|G|}+O\!\bigl(\tau(u)^n\bigr)\right),
\qquad
B_n(\ind;u)
=
\frac{\lambda_1(u)^n}{n}\left(1+O\!\bigl(\tau(u)^n\bigr)\right).
\]
Because \(\tau(u)<1\), for all sufficiently large \(n\) the second factor lies in \([1/2,3/2]\); in particular \(B_n(\ind;u)>0\), and
\[
\frac{1}{B_n(\ind;u)}
=
\frac{n}{\lambda_1(u)^n}\left(1+O\!\bigl(\tau(u)^n\bigr)\right).
\]
It follows that
\[
\mu_{n,u}(c)=\frac{1}{|G|}+O\!\bigl(\tau(u)^n\bigr)
\]
uniformly for \(c\in G\). Therefore,
\[
\|\mu_{n,u}-\mathrm{unif}_G\|_{\mathrm{TV}}
=
\frac12\sum_{c\in G}\left|\mu_{n,u}(c)-\frac{1}{|G|}\right|
=
O\!\bigl(\tau(u)^n\bigr),
\]
which proves the total-variation estimate.

Now set
\[
r_n(c):=|G|\mu_{n,u}(c)=1+O\!\bigl(\tau(u)^n\bigr),
\]
again uniformly in \(c\). Then
\[
\chi^2\!\bigl(\mu_{n,u}\,\|\,\mathrm{unif}_G\bigr)
=
\frac{1}{|G|}\sum_{c\in G}(r_n(c)-1)^2
=
O\!\bigl(\tau(u)^{2n}\bigr).
\]
For the relative entropy, Jensen's inequality for the concave function \(\log\) gives
\[
\begin{aligned}
D\!\bigl(\mu_{n,u}\,\|\,\mathrm{unif}_G\bigr)
&=
\sum_{c\in G}\mu_{n,u}(c)\log r_n(c)\\
&\le
\log\!\left(\sum_{c\in G}\mu_{n,u}(c)\,r_n(c)\right)\\
&=
\log\!\bigl(1+\chi^2(\mu_{n,u}\Vert \mathrm{unif}_G)\bigr).
\end{aligned}
\]
Since \(\log(1+x)\le x\) for \(x\ge 0\), the right-hand side is at most
\[
\chi^2\!\bigl(\mu_{n,u}\,\|\,\mathrm{unif}_G\bigr)
=
O\!\bigl(\tau(u)^{2n}\bigr).
\]
This proves the first three assertions.

If \(\lambda_1(u)>1\) and \(\textsf{SR}_{K,G}(u)\) holds, then \(\tau(u)\le \lambda_1(u)^{-1/2}<1\), and the preceding estimates immediately yield the square-root bounds.

For the converse statement, fix a subset \(A\subseteq G\) as in the hypothesis, and choose \(\Gamma_+\) with \(\Gamma(u)<\Gamma_+<\eta(u)\). By Theorem \ref{thm:kernel-chebotarev-second-main-term-witness},
\[
B_{n,A}(u)
=
\mathrm{unif}_G(A)\frac{\lambda_1(u)^n}{n}
+
\Re\!\bigl(c_A\mu_\star^n\bigr)\frac{1}{n}
+
O\!\left(\frac{\Gamma_+^n}{n}\right).
\]
The same theorem also yields a constant \(c_0(A)>0\) such that
\[
\bigl|\Re(c_A\mu_\star^n)\bigr|\ge 2c_0(A)\eta(u)^n
\]
for infinitely many \(n\). Since \(\Gamma_+<\eta(u)\), the error term is \(o(\eta(u)^n)\), so after possibly decreasing \(c_0(A)\) we obtain
\[
\left|B_{n,A}(u)-\mathrm{unif}_G(A)\frac{\lambda_1(u)^n}{n}\right|
\ge
c_0(A)\frac{\eta(u)^n}{n}
\]
for infinitely many \(n\). Since
\[
B_n(\ind;u)=\frac{\lambda_1(u)^n}{n}\bigl(1+o(1)\bigr),
\]
dividing by \(B_n(\ind;u)\) yields
\[
\bigl|\mu_{n,u}(A)-\mathrm{unif}_G(A)\bigr|
\ge
c(A)\,\frac{\eta(u)^n}{\lambda_1(u)^n}
\]
along infinitely many \(n\), after possibly decreasing the constant once more. If one had
\[
\|\mu_{n,u}-\mathrm{unif}_G\|_{\mathrm{TV}}
=
O\!\bigl(\lambda_1(u)^{-n/2}\bigr),
\]
then every subset discrepancy would satisfy the same bound because
\[
\bigl|\mu_{n,u}(A)-\mathrm{unif}_G(A)\bigr|
\le
\|\mu_{n,u}-\mathrm{unif}_G\|_{\mathrm{TV}}.
\]
This is impossible because
\[
\frac{\eta(u)^n}{\lambda_1(u)^n}
=
\left(\frac{\eta(u)}{\lambda_1(u)^{1/2}}\right)^n\lambda_1(u)^{-n/2},
\]
and the prefactor grows exponentially when \(\eta(u)>\lambda_1(u)^{1/2}\). Hence the square-root bound in total variation cannot hold in this regime.
\end{proof}

\endinput
